\chapter{Diseño experimental}\label{sec:diseno}

\section{Evaluaciones y unidades de análisis}
Se distinguen el benchmark estático, la cascada Top-$k$ de referencia y el desarrollo del rescate de ensamble. La Tabla~\ref{tab:cobertura} especifica sus unidades; las condiciones externas no son datasets independientes. El híbrido se evalúa con SVM y no hereda los resultados ResMLP del benchmark.

\begin{table}[!htbp]
\centering\small
\caption{Cobertura efectiva de las evaluaciones. Las condiciones de cada conjunto se agregan antes de la inferencia entre datasets.}
\label{tab:cobertura}
\begin{tabular}{@{}lrrr@{}}
\toprule
Conjunto & Imágenes/clases & Estático por clasif. & Cascada referencia\\
\midrule
DTD & 5.640/47 & 10 & 10\\
FMD & 1.000/10 & 15 & 5\\
CUReT & 5.612/61 & 2 & 2\\
Outex13 & 1.360/68 & 1 & 1\\
SoilOriginal & 1.140/7 & 15 & 5\\
KTH-TIPS-2b & 4.752/11 & 4 & 4\\
\midrule
Total de condiciones & & 47 & 27\\
\bottomrule
\end{tabular}

\vspace{0.5em}
\begin{minipage}{0.96\textwidth}\small
Piloto híbrido: entrenamiento oficial Outex, semilla 42. Réplica: FMD con semillas 42/123/2026 (800 imágenes de desarrollo por condición) y Outex con 123/2026 (680). Cada condición utiliza tres pliegues productores y tres internos. No se evalúan sus tests externos ni se añaden esas cinco condiciones al panel histórico de 27.
\end{minipage}
\end{table}

\section{Datos, particiones y procedencia}
Se utilizaron manifiestos para verificar etiquetas, orden de muestras, dimensiones, valores finitos y grupos de dependencia. Los bloques se alinearon con el mismo orden de imágenes dentro de cada conjunto.

DTD proporciona diez particiones oficiales \citep{cimpoi2014dtd}. En cada corrida estática se unieron entrenamiento y validación para seleccionar internamente; el test se reservó para evaluación externa de esa corrida. Se retiraron del entrenamiento los duplicados byte-idénticos que cruzaban roles. Esta separación no significa que los resultados oficiales nunca influyeran en decisiones posteriores del proyecto.

FMD contiene diez categorías de materiales \citep{sharan2013fmd}. Cada imagen tuvo un hash único y se utilizó como grupo. El benchmark repite cinco pliegues estratificados y agrupados con semillas 42, 123 y 2026; la cascada conserva los cinco de la semilla 42.

Outex-TC-00013 contiene 68 clases balanceadas y una partición oficial de 680 imágenes de entrenamiento y 680 de prueba, diez por clase a cada lado \citep{ojala2002outex}. Se verificaron archivos RGB legibles de $128\times128$, IDs continuos y ausencia de duplicados exactos. El piloto híbrido trabaja exclusivamente dentro de las 680 imágenes de entrenamiento.

CUReT utiliza la distribución gris recortada de 61 materiales y 92 condiciones por clase \citep{varma2003filter}. Se definieron mitades deterministas de 46 condiciones y ambas direcciones de entrenamiento/prueba. Las condiciones se agruparon para no cruzar una misma configuración de adquisición. El tamaño 46/46 no identifica el split histórico exacto de otros artículos; las dos direcciones son dependientes.

SoilOriginal procede de la versión 1 de \href{https://www.kaggle.com/datasets/ai4a-lab/comprehensive-soil-classification-datasets}{Comprehensive Soil Classification Datasets, de ai4a-lab}. Se admitió únicamente el árbol de imágenes originales del archivo auditado, excluyendo el conjunto aumentado. La deduplicación exacta por hash de archivo conserva 1.140 imágenes de siete clases y 962 grupos conservadores por familia de nombre, definidos dentro de cada clase. El manifiesto fija versión, rutas y hashes; no identifica un split oficial del proveedor. El benchmark utiliza cinco pliegues agrupados con tres semillas; la cascada, los cinco de la semilla 42, con 228 imágenes de prueba por pliegue.

KTH-TIPS-2b contiene 11 clases y cuatro muestras físicas por clase. Se conservaron cuatro particiones invertidas, con tres muestras físicas de entrenamiento y una de prueba. Las tablas estáticas principales conservan los controles archivados, incluido el Top-$k$ anterior, y sus contrastes se reportan como análisis de ese archivo. No se les atribuye la corrección interna por muestra de controles posteriores. La cascada construida con el rescate anterior tampoco queda reconciliada por sustituir una media. Por esa reserva, KTH se excluye del contraste conjunto de la cascada. El Anexo~\ref{app:procedencia} identifica las variantes y sus usos.

\section{Clasificadores y selección interna}
La SVM lineal utiliza \texttt{LinearSVC}, $C=1$, clases balanceadas, pérdida \emph{squared hinge}, regularización $L_2$, uno-contra-el-resto, tolerancia $10^{-4}$ y hasta 10.000 iteraciones. Se mantiene la misma receta al comparar concatenaciones en el piloto.

ResMLP designa una MLP residual sobre vectores de características, no una reproducción del modelo de imágenes homónimo \citep{touvron2021resmlp}. Proyecta a 256 unidades y aplica tres bloques residuales, normalización de capa, GELU y \emph{dropout} de 0,1. Se entrenó con entropía cruzada y AdamW, tasa $10^{-3}$, decaimiento $10^{-4}$, lotes de 256 y hasta 100 épocas. La parada controla pérdida media de entrenamiento, paciencia diez y disminución mínima $10^{-4}$, restaurando el mejor estado. No selecciona época con test ni con validación adicional.

La búsqueda estática utiliza hasta cuatro pliegues internos estratificados y agrupados, limitados por disponibilidad de grupos. Las variantes KTH se delimitan en la sección anterior. Las transformaciones aprendidas, incluida SVD de N-gramas, se ajustan en el entrenamiento pertinente. Las redes de extracción permanecen congeladas; no se comparan esquemas de ajuste fino.

\section{Piloto de rescate de ensamble}
Sea $D_{\mathrm{dev}}$ el entrenamiento oficial de Outex. Cada pliegue productor $h$ lo divide en entrenamiento $T_h$ y validación $V_h$, sin grupos compartidos. A su vez, $T_h$ se divide en tres pliegues internos. Se distinguen cuatro operaciones:
\begin{enumerate}
\item Construir el ranking, calcular CKA y generar los subconjuntos utilizando $T_h$.
\item Reunir puntuaciones OOF internas de las imágenes de $T_h$ y seleccionar base, umbral y rescate mediante la Ecuación~\ref{eq:policy-selection}.
\item Reajustar los modelos elegidos con $T_h$, estimar sus escalas y conservar el umbral numérico seleccionado.
\item Predecir las imágenes de $V_h$ sin modificar la política a partir de sus etiquetas.
\end{enumerate}
Al reunir las tres validaciones $V_h$ se obtienen las 680 predicciones OOF de desarrollo, una por imagen y política. Esas predicciones son distintas de las puntuaciones OOF internas que guiaron la selección.

La biblioteca de candidatos se construye a nivel de $T_h$; no se genera de nuevo dentro de cada pliegue interno. Las puntuaciones internas, por tanto, sirven como criterio de selección dentro de una búsqueda adaptada a $T_h$, no como evaluación independiente del procedimiento completo. La evaluación productora conserva separadas las imágenes de $V_h$. La semilla de partición productora es 142 y la interna, 42; el Anexo~\ref{app:procedencia} registra tamaños y multiplicidades.

El piloto reutiliza puntuaciones de validación productora de los miembros del control de diversidad previamente auditado. Se ajustaron las puntuaciones internas faltantes y las bases necesarias, comprobando igualdad de particiones y procedencia. Los parámetros del ensamble, presupuestos y cuantiles se fijaron desde la configuración de desarrollo; no son una preespecificación ciega de toda la investigación.

Los controles del mismo piloto incluyen mejor individual interno, Top-3/4/5, Completa-22 y ensamble de diversidad permanente. Top-3/4/5 son tamaños fijos; no equivalen al Top-$k$ del benchmark que hereda $k$ de GFS. Se añadió un ensamble estático seleccionado bajo cada presupuesto, definido a continuación; no se dispone de GFS concatenado reajustado para este protocolo.

\subsection{Ablación del mismo rescate con y sin compuerta}\label{sec:diseno-p0}
Se reconstruyeron cuatro vistas de cada una de las doce políticas a partir de puntuaciones y predicciones de validación productora guardadas: base sola, mismo rescate permanente, ruta forzada y ruta condicional. Se conservaron los miembros, multiplicidades, pesos, escalas, modelos y particiones; no se reentrenó ni se utilizó el test oficial. La base de cada vista es la elegida para esa política, no el mejor individual global.

El rescate permanente predice todas las imágenes con los miembros fijados y contabiliza únicamente la unión $U(E)$, sin extraer una base innecesaria. La ruta forzada conserva además la base y el sobrecosto de política, pero activa el mismo rescate en todas las imágenes. Ambas producen las mismas clases; la segunda tiene otra contabilidad de costo y llamadas. La vista condicional conserva la decisión archivada de la compuerta. Si el rescate es vacío, las cuatro vistas se reducen a la base; esos casos no aportan evidencia sobre enrutamiento.

El contraste entre rescate permanente y condicional aísla la activación para los modelos fijados en este archivo de desarrollo. No evalúa el mejor rescate estático posible: sus miembros fueron seleccionados conjuntamente con la compuerta. Las doce configuraciones reutilizan las mismas 680 imágenes y son dependientes. Se informan cambios predictivos, costos y llamadas, sin atribuirles significancia confirmatoria ni equivalencia poblacional.

\subsection{Control estático bajo presupuesto}\label{sec:diseno-estatico-presupuesto}
En cada entrenamiento productor se reúnen los 22 clasificadores individuales y los modelos únicos del ensamble de diversidad, con las puntuaciones OOF internas guardadas. Se enumeran todos los ensambles de uno a tres miembros únicos de ese pool, con igual peso y sin repetición. Para cada presupuesto 15, 30, 60 y 120 ms se elige el mayor macro-F1 interno cuya unión de bloques sea factible. Los desempates favorecen menor costo, menos modelos y orden léxico. No se añade una base ni una compuerta al costo del control permanente.

Las doce decisiones se guardaron antes de calcular sus predicciones de validación productora. Se reutilizaron nueve modelos disponibles y se reajustaron nueve con la misma receta SVM y el mismo entrenamiento productor. La enumeración es exhaustiva dentro de ese pool y límite de tres miembros, no sobre todas las concatenaciones o ponderaciones posibles. Se conservan los mismos pliegues y constantes del híbrido; el control no sustituye una ejecución GFS bajo presupuesto.

Este análisis se diseñó después de conocer el piloto y P0. Su regla se fijó antes de obtener las nuevas predicciones, pero no constituye una preinscripción ciega. Se conserva todo el barrido y no se usa el test oficial ni se reajusta el híbrido a partir del resultado del control.

\subsection{Réplica en FMD y semillas adicionales}\label{sec:diseno-replica}
La réplica conserva la biblioteca de 22 bloques, SVM, tres pliegues productores y tres internos, presupuestos 15/30/60/120 ms y cuantiles 0,10/0,25/0,50. Mantiene cinco sorteos de candidatos de tamaño cuatro, hasta tres miembros podados y el ensamble de quince votos con $\beta=8$ y $\lambda=0{,}25$. No se modificaron estos parámetros a partir de sus resultados. Las cinco condiciones y los comandos se fijaron antes de ejecutar la cola, después de conocer el piloto; no constituyen una preinscripción ciega de la investigación.

En FMD se reserva el pliegue externo 0 de la partición estratificada y agrupada de cada semilla 42, 123 y 2026: 800 imágenes forman $D_{\mathrm{dev}}$ y las 200 restantes quedan excluidas. En Outex se mantiene el entrenamiento oficial de 680 imágenes y se añaden las semillas 123 y 2026, sin evaluar las 680 imágenes de test. El piloto Outex/42 se conserva por separado. Las semillas alteran particiones, sorteos y transformaciones; no aíslan solamente el efecto de la inicialización. La partición productora usa semilla $s+100$ y la interna, $s$.

Cada condición evalúa referencias individuales, Top-3/4/5, Completa-22, ensambles permanentes, las tres reglas híbridas y el rival estático bajo presupuesto. La vista P0 vuelve a aplicar permanentemente exactamente el rescate seleccionado, sin reoptimizarlo. La cobertura se audita por imagen, grupo y transformación, manteniendo el test externo excluido. El perfil de costos es el mismo del piloto Outex, no una medición nueva de extracción en FMD. Las reglas completas y el detalle por semilla se conservan en el Anexo~\ref{sec:anexo-replica}.

\section{Métricas y agregación}
Macro-F1 es la métrica primaria. Para $C$ clases,
\begin{equation}
\operatorname{macro\mbox{-}F1}=\frac1C\sum_{c=1}^{C}
\frac{2TP_c}{2TP_c+FP_c+FN_c}.
\end{equation}
Se asigna cero cuando el denominador no permite definir el cociente. La exactitud es la proporción de predicciones correctas. En el benchmark se promedian métricas entre condiciones externas de cada dataset; en cada condición del híbrido se calcula macro-F1 sobre las predicciones OOF productoras reunidas, una por imagen y política (680 Outex, 800 FMD). Las medias de la réplica promedian después esas métricas entre semillas dentro de cada dataset; no reúnen varias predicciones de una misma imagen como observaciones independientes.

El costo proyectado del piloto se promedia por imagen, ponderando tamaños de pliegue; no se usa una media no ponderada de costos de pliegues de distinto tamaño. Las tablas principales muestran medias sin formato mixto de desviaciones. El Anexo~\ref{app:resultados-adicionales} conserva desviaciones descriptivas entre condiciones y el detalle del barrido; un único split Outex no permite estimar desviación entre splits.

\section{Medición física de la cascada de referencia}
Se midió una condición por dataset en RTX 4060 de 8 GB: DTD/Outex oficial 1, CUReT A a B, KTH oficial 1 invertido y FMD/Soil semilla 42, pliegue 0. Se seleccionó una imagen por clase, con semilla 42, para 204 imágenes y tres repeticiones pareadas de orden aleatorizado: 612 parejas, no 612 imágenes independientes.

Se utilizó lote unitario, modelos residentes e imágenes decodificadas, con sincronización CUDA. El cronómetro incluyó extracción y clasificación, excluyendo inicialización y disco. Se contrastaron representaciones y predicciones con el caché. Las discrepancias de CUReT/Soil limitan la fidelidad; esta medición no corresponde al híbrido nuevo.

\section{Panel físico online completo (94 condiciones) y comparadores del estado del arte}\label{sec:diseno-panel-online}
Para superar la cobertura muestral limitada de los pilotos físicos iniciales y comparar directamente tanto calidad predictiva como tiempo físico de inferencia por imagen frente al estado del arte reciente de concatenación, se ejecutó un panel online completo en lote unitario (\texttt{batch=1}) sobre las 47 particiones externas de los seis conjuntos de datos (FMD: 15, DTD: 10, SoilOriginal: 15, KTH-TIPS-2b: 4, CUReT: 2 y Outex-13: 1):
\begin{enumerate}
\item \textbf{Panel local online (94 condiciones, 362.640 llamadas primarias):} En cada una de las 47 particiones externas y para ambos clasificadores (\texttt{LinearSVC} y \texttt{ResMLP}, totalizando 94 condiciones únicas y 470 combinaciones condición--método), se evaluaron cinco políticas locales sobre la totalidad de las imágenes de prueba: \texttt{Adaptive-selection-margin} (adquisición condicional con salida temprana por margen en la Etapa 1 y rescate en la Etapa 2), \texttt{Top-4} (concatenación estática de los cuatro mejores bloques individuales), \texttt{Ens-diversidad} (selección estática de cuatro bloques guiada por relevancia y diversidad CKA), \texttt{Ens-tombola} (selección estática ponderada de cuatro bloques) y \texttt{Completa-22} (concatenación de los 22 bloques). Cada llamada individual midió extremo a extremo la extracción y clasificación con sincronización CUDA (\texttt{perf\_counter\_ns}), derivando el macro-F1 y la latencia física (\texttt{ms}) exactamente de las mismas predicciones individuales, sin reutilizar cachés en línea.
\item \textbf{Validación de fidelidad y adaptación de cuatro métodos externos de concatenación:} Se seleccionaron cuatro artículos representativos del estado del arte que emplean concatenación de características profundas y/o artesanales:
\begin{itemize}
\item \textbf{Neshov et al.\ (2025) \citep{neshov2025aggregation}:} Agregación multicapa de los 12 bloques Transformer de \texttt{BEiT-v2} (\texttt{beitv2\_base\_patch16\_224.in1k\_ft\_in22k\_in1k}) mediante \emph{Global Average Pooling} concatenado ($12\times 768 = 9.216$ dimensiones) y clasificador MLP con normalización por lotes (\texttt{BN $\to$ Linear(9216,64) $\to$ BN $\to$ ReLU $\to$ Dropout(0,2) $\to$ Linear(64,$C$)}, Adam $10^{-3}$, 100 épocas). Se verificó primero su protocolo publicado en FMD ($10\times 10$-fold CV, 100 condiciones y 10.000 llamadas) y luego se evaluó su adaptación con extracción individual \texttt{batch=1} en las 47 condiciones del panel común (36.264 llamadas).
\item \textbf{Simon y Uma (2022) --- \emph{DeepLumina} \citep{simon2022deeplumina}:} Conversión de la imagen RGB al canal de luminancia $Y = 0{,}299R + 0{,}587G + 0{,}114B$ del espacio YIQ, extracción independiente de la capa \texttt{fc1000} de \texttt{ResNet-101} sobre la entrada RGB (1.000D) y sobre el canal $Y$ replicado en tres canales (1.000D) para formar un vector concatenado de 2.000 dimensiones, clasificado mediante un diseño multiclase ECOC-SVM \emph{sparse-random} ($\lceil 15\log_2(C)\rceil$ columnas iniciales evaluadas sobre 10.000 matrices ternarias candidatas maximizando la distancia mínima entre pares de clases y decodificación \emph{loss-weighted hinge}).
\item \textbf{Gupta et al.\ (2026) --- \emph{HyTexNet} \citep{gupta2026hytexnet}:} Concatenación de la representación profunda de \texttt{DenseNet-121} (1.024D) con el descriptor local \emph{HyTex} de 514 dimensiones (histogramas de signo, magnitud umbralizada en el percentil 75 empírico $T_{75}$ y media global central, totalizando 1.538 dimensiones) seguida de una cabeza MLP (\texttt{1538 $\to$ 256 $\to$ ReLU $\to$ Dropout(0,5) $\to$ $C$}). Se validaron mediante pruebas unitarias los cinco ejemplos numéricos paso a paso de la Tabla 1 del artículo original, el protocolo de un tercio de entrenamiento en KTH-TIPS-2b (3.168 llamadas) y las 47 condiciones del panel común (36.264 llamadas).
\item \textbf{Chang et al.\ (2024) --- \emph{6LBP+dCellNN} \citep{chang2024cellnn}:} Fusión de mapas de contorno extraídos con el checkpoint oficial de los autores (\texttt{cellnn\_model.ph}, commit \texttt{299502c} del repositorio público \texttt{UESTC-nnLab/LBP\_Contour}) y patrones binarios locales en seis escalas $(P,R)\in\{(8,1),(16,2),(24,3),(8,3),(12,5),(16,7)\}$ concatenados en 768 dimensiones con clasificador MLP de 400 unidades ocultas. Se verificó paridad numérica exacta ($\text{error máximo} = 0{,}0$) frente al código original del autor antes de evaluar tanto el protocolo de un tercio en KTH-TIPS-2b como las 47 condiciones del panel común.
\end{itemize}
\end{enumerate}

\section{Análisis estadístico}\label{sec:analisis-estadistico}
Los análisis siguen la comparación no paramétrica entre datasets \citep{demsar2006statistical,derrac2011practical,garcia2008extension}, con cálculos originales en ControlTest/MultipleTest de Java y comprobaciones locales posteriores. En estos contrastes de calidad predictiva, la unidad de inferencia es el dataset, después de agregar sus condiciones; las imágenes y pliegues de un mismo conjunto no incrementan ese número de unidades.

Para el benchmark estático se utiliza una media por dataset y clasificador: $N=6$ bloques, con las reservas KTH descritas. Friedman es el contraste global principal de ese análisis histórico; sus post-hoc de diferencias de rangos aplican Holm a los diez pares por clasificador y no son Wilcoxon exactos. Iman--Davenport es una transformación del estadístico de Friedman, no una réplica independiente de evidencia. Friedman alineado y Quade son comprobaciones complementarias con sensibilidades distintas; su desacuerdo se informa sin escoger retrospectivamente la prueba más favorable.

Para examinar la reserva de KTH se repiten esas cuatro pruebas y las dos familias de diez pares excluyendo ese conjunto ($N=5$). Se mantienen intactos los resultados de los otros datasets y se reproducen primero los contrastes Java de $N=6$. Es una sensibilidad del mismo archivo, no una corrección de KTH ni una confirmación independiente; sus valores ajustados pertenecen a familias distintas de las originales.

Para la cascada de referencia se excluye KTH y se utilizan cinco datasets, correspondientes a 23 condiciones. La familia de tres contrastes pareados compara la cascada seleccionada con cascada ResNet, Top-$k$ de origen y Completa-22, mediante Wilcoxon y Holm. Friedman y sus post-hoc examinan rangos entre métodos; Wilcoxon examina diferencias pareadas respecto de cero bajo sus supuestos, no una hipótesis idéntica a la de Friedman.

El no rechazo no demuestra equivalencia ni no inferioridad. Las particiones y repeticiones compartidas no se tratan como datasets nuevos. No se realizan contrastes confirmatorios del punto favorable del piloto Outex entre presupuestos explorados. La réplica del híbrido se informa descriptivamente: tiene dos datasets y semillas dependientes, no cinco dominios independientes para justificar una nueva batería de contrastes poblacionales.
